\subsection{分块乘积}

第 4 号的定义可以如下推广。设 E 是一个（右或左）A-模，是一个子模族 \((\mathbf{E}_i)_{i\in I}\) 的直和。对所有 \(x\in \mathbf{E}\) ，设 \(x = \sum_{i\in I}x_i\) ，其中 \(x_{i}\in \mathbf{E}_{i}\) （对所有 \(i\in I\) ）；我们称元素属于 E 的列矩阵 \(M(x) = (x_{i})_{i\in I}\) 为 \(x\) 关于 E 的直和分解 \((\mathbf{E}_i)_{i\in I}\) 的矩阵。

设 F 是另一个 A-模（E 与 F 同为右 A-模或同为左 A-模），并假设 F 是子模族 \((\mathbf{F}_{k})_{k\in\mathbb{K}}\) 的直和。对所有 \(u\in\mathrm{Hom}(\mathbf{E},\mathbf{F})\) 与所有 \(x_{i}\in\mathbf{E}_{i}\) ，设 \(u(x_{i})=\sum_{k}u_{ki}(x_{i})\) ，其中 \(u_{ki}(x_{i})\in\mathbf{F}_{k}\) （对所有 \(k\in K\) ）；于是 \(u_{ki}\in\mathrm{Hom}(\mathbf{E}_{i},\mathbf{F}_{k})\) ；我们称矩阵 \(M(u)=(u_{ki})_{(k,i)\in K\times I}\) —— (K, I) 型、元素属于集合 H（即诸 \(\mathrm{Hom}(\mathbf{E}_{i},\mathbf{F}_{k})\) 之和）—— 为 u 关于 E 与 F 的直和分解 \((\mathbf{E}_{i})\) 与 \((\mathbf{F}_{k})\) 的矩阵。

有了这些定义，显然若 \(u, v\) 是 \(E\) 到 \(F\) 的两个 A-线性映射，则对关于相同直和分解的矩阵，有

\[
M (u + v) = M (u) + M (v), \quad M (\gamma u) = \gamma M (u)\tag{21}
\]

（对 A 的中心的每个元素 \(\gamma\) ）（第 2 号，注记）。

此外，诸 \(u_{ki}\) 的定义表明，若 K 有限，则我们可写出

\[
M (u (x)) = M (u). M (x)\tag{22}
\]

其中 \(M(u(x))\) 是 \(u(x)\) 关于分解 \((\mathbf{F}_k)\) 的矩阵，(22) 右边的乘积是对映射 \((t, z) \mapsto t(z)\) （从 \(\operatorname{Hom}(\mathbf{E}_i, \mathbf{F}_k) \times \mathbf{E}_i\) 到 \(\mathbf{F}_k\) ）计算的（第 2 号，注记）。

设 G 是第三个 A-模，是一个子模族 \((\mathbf{G}_{l})_{l\in L}\) 的直和，于是对每个 A-线性映射 \(v:F\to G\) 对应着一个矩阵 \(M(v)=(v_{lk})\) （关于分解 \((\mathbf{F}_{k})\) 与 \((\mathbf{G}_{l})\) ）。若 I、K 与 L 都有限，则

\[
M (v \circ u) = M (v). M (u)\tag{23}
\]

其中左端是矩阵 \((w_{li})\) —— \(w = v \circ u\) 关于分解 \((\mathbf{E}_i)\) 与 \((\mathbf{G}_l)\) 的矩阵——，且左端的乘积是对映射 \((t, s) \mapsto t \circ s\) （从 \(\operatorname{Hom}(\mathbf{F}_k, \mathbf{G}_l) \times \operatorname{Hom}(\mathbf{E}_i, \mathbf{F}_k)\) 到 \(\operatorname{Hom}(\mathbf{E}_i, \mathbf{G}_l)\) ）计算的（第 2 号，注记）。这正是 §1 第 8 号的公式 (32) 用矩阵写出的形式。

最后，若假设 I 与 K 有限，则 \(E^{*}\) （相应地 \(F^{*}\) ）典范地等同于诸模 \(E_{i}^{*}\) （相应地 \(F_{k}^{*}\) ）的直和（ \(\S\) 2，第 6 号，命题 10）。于是立即可验证， \(^{t}u\) 关于分解 \((\mathbf{F}_{k}^{*})\) 与 \((\mathbf{E}_{i}^{*})\) 的矩阵正是 \((^{t}u_{kt})_{(k,t)\in\mathbb{K}\times\mathbb{I}}\) 。

现在设 I 与 K 有限，并且进一步设每个 \(E_{i}\) （相应地 \(F_{k}\) ）都具有一个有限基。这等价于说 E（相应地 F）具有一个基 \((e_{r})_{r\in\mathbb{R}}\) （相应地 \((f_{s})_{s\in\mathbb{S}}\) ），且 R（相应地 S）具有一个划分 \((\mathbf{R}_{i})_{i\in\mathbb{I}}\) （相应地 \((\mathbf{S}_{k})_{k\in\mathbb{K}}\) ），使得对所有 \(i\in I\) （相应地 \(k\in K\) ）， \((e_{r})_{r\in\mathbb{R}_{i}}\) 是 \(E_{i}\) 的基（相应地 \((f_{s})_{s\in\mathbb{S}_{k}}\) 是 \(F_{k}\) 的基）。于是，若 \(X=M(u)\) 是 u 关于基 \((e_{r})_{r\in\mathbb{R}}\) 与 \((f_{s})_{s\in\mathbb{S}}\) 的矩阵，则矩阵 \(X_{ki}=M(u_{ki})\) （关于基 \((e_{r})_{r\in\mathbb{R}_{i}}\) 与 \((f_{s})_{s\in\mathbb{S}_{k}}\) ）正是 X 的这样一个子矩阵：它由删去指标为 \(s\notin S_{k}\) 的行与指标为 \(r\notin R_{i}\) 的列而得到。这样我们便定义了一个一一对应

\[
X \mapsto (X _ {k i}) _ {(k, i) \in \mathbf {K} \times \mathbf {I}}\tag{24}
\]

介于元素属于 A 的 (S, R) 型矩阵的集合与矩阵的矩阵 \((\mathbf{X}_{ki})_{(k,i)\in\mathbf{K}\times\mathbf{I}}\) （ \(K\times I\) 型，其中每个 \(X_{ki}\) 是 \((\mathbf{S}_{k},\mathbf{R}_{i})\) 型的、元素属于 A 的矩阵）的集合之间。进一步设 G 具有一个有限基 \((g_{t})_{t\in\mathbf{T}}\) ，且 \(\mathbf{T}=(\mathbf{T}_{l})_{l\in\mathbf{L}}\) 是 T 的一个划分，使得对每个 \(l\in L\) ， \((g_{t})_{t\in\mathbf{T}_{l}}\) 是 \(G_{l}\) 的基；令 \(Y=M(v)\) 是 v 关于基 \((f_{s})_{s\in\mathbf{S}}\) 与 \((g_{t})_{t\in\mathbf{T}}\) 的矩阵， \(Y_{lk}=M(v_{lk})\) 是 \(v_{lk}\) 关于基 \((f_{s})_{s\in\mathbf{S}_{k}}\) 与 \((g_{t})_{t\in\mathbf{T}_{l}}\) 的矩阵，Z=M(w) 是 w=v∘u 关于基 \((e_{r})_{r\in\mathbf{R}}\) 与 \((g_{t})_{t\in\mathbf{T}}\) 的矩阵，而 \(Z_{li}=M(w_{li})\) 是 \(w_{li}\) 关于基 \((e_{r})_{r\in\mathbf{R}_{i}}\) 与 \((g_{t})_{t\in\mathbf{T}_{l}}\) 的矩阵；于是由 (23)，Z=YX 的子矩阵 \(Z_{li}\) 由下式给出

\[
Z _ {l i} = \sum_ {k} Y _ {l k} X _ {k i}\tag{25}
\]

换言之，当所涉及的所有乘积都有定义时，一一对应 (24) 把乘积变为乘积（矩阵的矩阵的乘积按第 2 号注记的意义定义）；当 Z=YX 的子矩阵 \(Z_{li}\) 如此计算时，就称该乘积是“按分块”进行的。

这个名称来自如下事实：当 \(I = [1, p]\) 且 \(K = [1, q]\) 时，表示矩阵 X 的表格被设想成划分成一些“块”，构成一个“矩阵的阵列”

\[
\left( \begin{array}{c c c c} X _ {1 1} & X _ {1 2} & \ldots & X _ {1 p} \\ X _ {2 1} & X _ {2 2} & \ldots & X _ {2 p} \\ \cdot & \cdot & \cdot & \cdot \\ X _ {q 1} & X _ {q 2} & \ldots & X _ {q p} \end{array} \right)
\]

当 p 与 q 是足够小的特定整数而使此做法可行时，它被看作表示 X 的符号。

